\section{Memory Management}

Understanding how memory works in C++ is crucial for writing efficient, bug-free programs. This section covers the fundamentals of memory allocation, the difference between stack and heap, and the principles that guide resource management in C++.

\subsection{Stack vs. Heap Memory}

C++ programs use two primary regions of memory for storing data:

\subsubsection{The Stack}

The \textbf{stack} is a region of memory that stores local variables and function call information. It operates in a last-in, first-out (LIFO) manner.

\begin{itemize}
	\item \textbf{Automatic management}: Memory is allocated when a variable comes into scope and deallocated when it goes out of scope
	\item \textbf{Fast allocation}: Stack allocation is extremely fast (just moving a pointer)
	\item \textbf{Limited size}: Typically 1-8 MB depending on the system
	\item \textbf{Contiguous}: Variables are stored in contiguous memory
\end{itemize}

\begin{lstlisting}
	void exampleFunction() {
		int x = 10;           // Allocated on the stack
		double y = 3.14;      // Allocated on the stack
		int arr[100];         // Array allocated on the stack
		
		// When this function returns, x, y, and arr are automatically deallocated
	}
\end{lstlisting}

\subsubsection{The Heap}

The \textbf{heap} (also called the \textbf{free store}) is a region of memory used for dynamic allocation—memory that you explicitly request at runtime.

\begin{itemize}
	\item \textbf{Manual management}: You must explicitly allocate and deallocate
	\item \textbf{Slower allocation}: Heap allocation involves more overhead
	\item \textbf{Large size}: Limited only by available system memory
	\item \textbf{Non-contiguous}: Allocations may be scattered throughout memory
\end{itemize}

\begin{lstlisting}
	void exampleFunction() {
		int* ptr = new int(42);      // Allocated on the heap
		int* arr = new int[1000];    // Array allocated on the heap
		
		// Must manually deallocate!
		delete ptr;
		delete[] arr;
	}
\end{lstlisting}

\subsection{Dynamic Memory Allocation}

\subsubsection{The \texttt{new} Operator}

The \texttt{new} operator allocates memory on the heap and returns a pointer to it:

\begin{lstlisting}
	// Allocate a single value
	int* p1 = new int;          // Uninitialized
	int* p2 = new int(42);      // Initialized to 42
	int* p3 = new int{42};      // C++11 uniform initialization
	
	// Allocate an array
	int* arr = new int[10];     // Array of 10 integers (uninitialized)
	int* arr2 = new int[10]();  // Array of 10 integers (zero-initialized)
	int* arr3 = new int[5]{1, 2, 3, 4, 5};  // C++11 initializer list
\end{lstlisting}

\subsubsection{The \texttt{delete} Operator}

The \texttt{delete} operator deallocates memory previously allocated with \texttt{new}:

\begin{lstlisting}
	int* p = new int(42);
	delete p;       // Deallocate single value
	
	int* arr = new int[10];
	delete[] arr;   // Deallocate array - note the []
\end{lstlisting}

\begin{remark}
	\textbf{Critical}: Always use \texttt{delete[]} for arrays allocated with \texttt{new[]}. Using \texttt{delete} instead of \texttt{delete[]} on an array causes undefined behavior.
\end{remark}

\subsection{Memory Problems}

Dynamic memory management is error-prone. Here are common problems:

\subsubsection{Memory Leaks}

A \textbf{memory leak} occurs when allocated memory is never deallocated:

\begin{lstlisting}
	void leakyFunction() {
		int* p = new int(42);
		// Oops! We never delete p
		// The memory is leaked when this function returns
	}
	
	void anotherLeak() {
		int* p = new int(42);
		p = new int(100);  // Original memory is leaked!
		delete p;          // Only deletes the second allocation
	}
\end{lstlisting}

\subsubsection{Dangling Pointers}

A \textbf{dangling pointer} points to memory that has been deallocated:

\begin{lstlisting}
	int* p = new int(42);
	delete p;
	// p is now a dangling pointer
	*p = 100;  // Undefined behavior! Accessing freed memory
\end{lstlisting}

Best practice: Set pointers to \texttt{nullptr} after deleting:

\begin{lstlisting}
	int* p = new int(42);
	delete p;
	p = nullptr;  // Now it's clear the pointer is invalid
\end{lstlisting}

\subsubsection{Double Delete}

Deleting the same memory twice causes undefined behavior:

\begin{lstlisting}
	int* p = new int(42);
	delete p;
	delete p;  // Undefined behavior! Double delete
\end{lstlisting}

\subsubsection{Mismatched new/delete}

Using the wrong form of delete:

\begin{lstlisting}
	int* arr = new int[10];
	delete arr;    // WRONG! Should be delete[]
	
	int* p = new int(42);
	delete[] p;    // WRONG! Should be delete
\end{lstlisting}

\subsection{RAII: Resource Acquisition Is Initialization}

\textbf{RAII} is a fundamental C++ idiom that ties resource management to object lifetime. The principle is simple:

\begin{itemize}
	\item Acquire resources in a constructor
	\item Release resources in the destructor
	\item Let scope-based destruction handle cleanup automatically
\end{itemize}

\begin{example}[RAII Wrapper Class]{\leavevmode}
	
	\begin{lstlisting}
		class IntArray {
			private:
			int* data;
			size_t size;
			
			public:
			// Constructor acquires the resource
			IntArray(size_t n) : data(new int[n]), size(n) {}
			
			// Destructor releases the resource
			~IntArray() {
				delete[] data;
			}
			
			// Access elements
			int& operator[](size_t i) { return data[i]; }
		};
		
		void useArray() {
			IntArray arr(100);  // Memory allocated
			arr[0] = 42;
			// When arr goes out of scope, destructor automatically frees memory
		}  // No leak possible!
	\end{lstlisting}
\end{example}

\begin{remark}
	RAII is why C++ can be both low-level (manual memory control) and safe (automatic cleanup). Smart pointers, covered in a later section, are the standard library's RAII wrappers for dynamic memory.
\end{remark}

\subsection{When to Use Dynamic Allocation}

Use stack allocation (local variables) by default. Use heap allocation when:

\begin{enumerate}
	\item \textbf{Size is unknown at compile time}: You don't know how much memory you need until runtime
	\item \textbf{Size is too large for the stack}: Large arrays or data structures
	\item \textbf{Lifetime must extend beyond the current scope}: Data needs to outlive the function that created it
	\item \textbf{Polymorphism}: Storing derived class objects through base class pointers
\end{enumerate}

\begin{remark}
	In modern C++, prefer smart pointers (\texttt{std::unique\_ptr}, \texttt{std::shared\_ptr}) and standard containers (\texttt{std::vector}, \texttt{std::string}) over raw \texttt{new}/\texttt{delete}. These handle memory management automatically while following RAII principles.
\end{remark}